\section{Selmer signatures on the actual representation family}
\label{sec:selmer-signature-v052}

\status{SSC1--SSC16 prove the full original Selmer-signature comparison,
including integral trace conjugacy, arbitrary local injective changes of
coefficient DVR, the original crystalline local conditions, infinite
lengths and the tame-conductor term. SSC17 removes the second of the two
remaining suppliers in DFD15. The coherent Wach functor and comparison
remain an explicit input to the family setup; the actual common
\texttt{Concl\allowbreak F} producer is not claimed here.}

Fix the original $D:\texttt{D7}$ and write $G=G_E$, $O_i$,
$K_i=\operatorname{Frac}(O_i)$, $\ell_i$, $X_i$ and
$Z_i^{\mathrm{good}}$ as in DFD. For a coefficient domain $O$ put
\[
 A_O^r=\operatorname{Frac}(O)^r/O^r.
\]
It has the discrete topology. An integral matrix representation acts
on this quotient by the same matrices. We use continuous cochains;
\(H^1(G,A)\) always means cocycles modulo coboundaries in that sense.
For a DVR $O$ and natural number $\ell$, let
\[
 n_O(M)=
 \begin{cases}
 \infty,&\operatorname{length}_O M=\infty,\\
 \displaystyle\frac{\operatorname{length}_O M}
 {\operatorname{length}_O(O/\ell O)},
       &\operatorname{length}_O M<\infty.
 \end{cases}                                                    \tag{SSC.1}
\]
The finite denominator in this notation uses the source's
\(\mathrm{toNat}\) convention if a module has infinite length, and
division by zero is zero in $\mathbb Q$. At every actual location
$\ell_i\ne0$ is a nonunit and the denominator is positive and finite.
The extended nonnegative version \texttt{lenN} uses division in
$\mathbb R_{\ge0}\cup\{\infty\}$; we will distinguish these conventions
when comparing them.

For a signature $s$, $T_s$ denotes the source's \texttt{latT}:
the representation $\rho_s$ twisted by the inverse of its specified
first character. For an injective coefficient map $f$ and character
$\psi$, write $T_{s,f,\psi}=f(T_s)\psi^{-1}$. The local condition
\texttt{extCond} is defined on cocycles: take the coefficient-ring
span of pushes of the original allowed cocycles, intersect with the
cocycles for $T_{s,f,\psi}$, and then pass to cohomology. Keeping this
order is necessary when $\psi$ need not be trivial locally.
The Selmer module \texttt{selT} uses this condition above $\ell$,
no condition at the specified bad and auxiliary places, and the
unramified condition elsewhere. All places, embeddings, characters,
and bad sets below are those supplied by $D$.

\subsection{Integral representations and the local conditions}

% BEGIN SSC1
\paragraph{SSC1: trace equality gives integral conjugacy.}
\label{claim:SSC1}
Let $O$ be any commutative local ring, $G$ any group, and
$\rho,\rho':G\to\operatorname{GL}_a(O)$. If the reduction of
$\rho$ is absolutely irreducible and the two traces agree on $G$,
then some $U\in\operatorname{GL}_a(O)$ satisfies
\[
                    \rho'(g)U=U\rho(g)\quad(g\in G).             \tag{SSC.2}
\]
The case $a=0$ is immediate. For $a>0$, absolute irreducibility
and the finite-dimensional density theorem imply that the residual
matrices span $M_a(k)$, where $k$ is the residue field. One may
check the use of this standard theorem after an algebraic closure:
Burnside's theorem gives the full matrix algebra there, and dimension
descends under field extension. Select $a^2$ original matrices whose
reductions form a basis. Their coordinate determinant is a unit, so
they are already a basis of $M_a(O)$. Thus the group-algebra map
$\phi:O[G]\to M_a(O)$ induced by $\rho$ is onto.

Here is a direct local-ring argument for the remaining assertion.
Write $\phi'$ for the other group-algebra map and choose an
$O$-linear right inverse $s$ to $\phi$. Set $\theta=\phi's$ and
\[
 T(X)_{ij}=\operatorname{tr}\bigl(X\theta(E_{ji})\bigr).
\]
For every $x\in O[G]$, trace equality applied to $xs(E_{ji})$
gives $T(\phi'(x))_{ij}=\phi(x)_{ij}$. Consequently
$T\theta=1$. The square matrix of $T$ has a right inverse,
so its determinant is a unit and $T$ is bijective. It follows that
$\theta\phi=\phi'$. Surjectivity of $\phi$ now implies that
$\theta$ is a unital algebra homomorphism and preserves traces.

Fix an index $o$. Since $\operatorname{tr}\theta(E_{oo})=1$,
some diagonal entry $u=\theta(E_{oo})_{yy}$ is a unit: all
nonunits of a local ring lie in its maximal ideal. Define
\[
 U_{xi}=\theta(E_{io})_{xy},\qquad
 V_{jx}=\theta(E_{oj})_{yx}.
\]
The matrix-unit multiplication rules give $VU=uI$. Taking
determinants proves that $U$ is invertible. Those same rules,
applied to $XE_{io}=\sum_k X_{ki}E_{ko}$, give
$\theta(X)U=UX$. Substitute $X=\rho(g)$ to obtain (SSC.2).
This proves the original DVR conjugacy theorem and, in fact, its
stronger local-ring version, without choosing a new integral lattice.
% END SSC1

% BEGIN SSC2
\paragraph{SSC2: the actual rank-$r$ summand has the triple's trace.}
\label{claim:SSC2}
Retain the original \texttt{Cheb} premise. Fix $i,\pi$, an
\texttt{IsRhoPt} representation $\rho_i(\pi)$, one list $S_g$
satisfying \texttt{HasParamX}, and a signature $s$ satisfying
\texttt{IsTriple} for that same list. Then
\[
 \operatorname{tr}\rho_i(\pi)(g)
   =f_i\bigl(\operatorname{tr}\rho_s(g)\bigr),\qquad g\in G,
 \quad f_i=D.\mathrm{ntw}.\mathrm{twist}_2(i).                  \tag{SSC.3}
\]
Indeed the actual rational splitting in \texttt{IsRhoPt} gives
\[
 \operatorname{tr}R_\pi(g)
    =\operatorname{tr}\rho_i(\pi)(g)
       +\sum_{c=1}^{t}\chi_{i,c}(g).                           \tag{SSC.4}
\]
This follows by conjugating to the two diagonal blocks; it makes
no integral splitting assertion. At every permitted split Frobenius,
the Hecke polynomial identifies the trace on the left with the
normalized parameter of $S_g$ together with its specified character
block. The triple identifies the trace of $\rho_s$ with the parameter
of $S_g$ alone. In the source normalization the character block is
\[
 \sum_{c=0}^{t-1}\iota_{QC}(\bar\chi_{c+1}(\sigma_v))
   =\operatorname{sumPart}(\mu\eta,t;v)
          (Nv)^{(1-n_V)/2}.                                   \tag{SSC.5}
\]
To check this equality, insert
$\bar\chi_{c+1}=\bar\chi_1\bar\varepsilon^{-c}$ in each term.
The stated Frobenius values are $\bar\varepsilon(\sigma_v)=Nv$
and, after the fixed complex embedding,
$\bar\chi_1(\sigma_v)=\operatorname{sumEig}(\mu\eta;v)
(Nv)^{(t-1)/2+(1-n_V)/2}$. Add the real exponents in each
of the $t$ summands. The norm $Nv$ is a positive integer,
so the ordinary real-power addition law is valid. The same
normalizing factor occurs in both parameter identities; (SSC.5)
cancels the character summand of (SSC.4), giving (SSC.3) on the
permitted Frobenius set.

Both traces in (SSC.3) are continuous in the original $\ell$-adic
coefficient field. In particular $f_i$ is continuous: the inverse
image of the target maximal ideal is a nonzero prime of the source
DVR, hence its maximal ideal, and powers of that ideal map into
the corresponding powers in the target. The original embeddings
into the $\ell$-adic coefficient field are continuous. DFD3, with
the same finite avoidance and split-place requirements, extends the
identity to $G$. No continuity of an embedding into $\mathbb C$
has been used.

If the residual representation is absolutely irreducible, SSC1
gives an integral conjugacy with $f_i(\rho_s)$. The triple identifies
the first character with the product of the pushed triple character
and $\psi_L(i)$. Multiplying (SSC.2) by these scalar inverses gives
\[
 W_{\pi}(g)U^{-1}
     =U^{-1}T_{s,f_i,\psi_L(i)}(g).                            \tag{SSC.6}
\]
Here $W_\pi$ is exactly \texttt{WpO}; inversion of $U$ accounts for
the source's direction of the intertwiner. This proves the full
\texttt{twistedRho\_conj\_latTpsi} implication on the original
integral representations.
% END SSC2

% BEGIN SSC3
\paragraph{SSC3: the actual local bridges transport crystalline cocycles.}
\label{claim:SSC3}
Two embeddings $j_1,j_2:\Omega\to\overline{\mathbb Q}_\ell$
inducing the same valuation ring are related by an automorphism
$\sigma$ over $\mathbb Q_\ell$. We include the argument needed
to compare the source's independently specified bridges. Equality
of the closed unit balls gives equality of the open unit balls by
applying it to $x^{-1}$. For a finite tuple $x_a\in\Omega$,
the condition
$\|\sum c_a j_t(x_a)\|<1$ is clopen as a condition on
$(c_a)\in\mathbb Q_\ell^n$. The two conditions agree on rational
tuples, so density of $\mathbb Q^n$ makes them agree everywhere.
If a linear combination is zero for $j_1$ and nonzero for $j_2$,
multiply all coefficients by a scalar of sufficiently large norm
to contradict this agreement. Hence the two maps
$\mathbb Q_\ell\otimes_{\mathbb Z}\Omega\to
\overline{\mathbb Q}_\ell$ have the same kernel. The induced
isomorphism of their image fields extends to an automorphism of
the algebraic closure, as asserted.

When $\Omega$ is algebraically closed, its image is dense in
$\overline{\mathbb Q}_\ell$. For completeness, approximate the
coefficients of a monic polynomial of an algebraic $\ell$-adic
element by rational coefficients. Factoring the approximating
polynomial shows that one root is arbitrarily close to the chosen
element: otherwise its value there, a product of root differences,
could not tend to zero. Such roots belong to the image of
$\Omega$. All $\mathbb Q_\ell$-automorphisms are isometries, so
agreement on $j(\Omega)$ implies agreement everywhere. The two
decomposition-group identifications in the bridges therefore obey
the exact conjugacy
$g_2=\sigma g_1\sigma^{-1}$, rather than an unspecified
abstract group isomorphism.

Here is the cochain transport. A crystalline cocycle $c$ has a
witness $x\in V\otimes_{\mathbb Q_\ell}B$ with
$c(g)\otimes1=gx-x$, where $B$ is the period ring of the supplied
\texttt{LocPrime}. If $q:V_1\to V_2$ is a
$\mathbb Q_\ell$-linear map intertwining the representations
under the above group conjugacy, use
\[
 q\otimes\sigma_B:V_1\otimes B\longrightarrow V_2\otimes B.
\]
On pure tensors the conjugacy identity proves equivariance;
linearity proves it everywhere. Applying this map to $x$ proves
that $q(c(g_1))$ is crystalline, with the displayed image of $x$
as witness. Its continuity follows from finite-dimensional
continuity of $q$ and continuity of the group maps. This proves
the full original \texttt{crystallineCocycle\_transport\_conjugation}
with the actual supplied period-ring action.

The needed coefficient maps automatically respect
$\mathbb Z_\ell$. In a Noetherian local target with $\ell$ in its
maximal ideal, two homomorphisms from $\mathbb Z_\ell$ agree:
approximate $z$ by an integer modulo $\ell^n$; their values then
differ by an element of every power of the maximal ideal, whose
intersection is zero. This supplies the precise scalar compatibility
for the coefficient maps used below.
% END SSC3

% BEGIN SSC4
\paragraph{SSC4: integral coefficient extension and conjugacy preserve the local conditions.}
\label{claim:SSC4}
Let $O_0\hookrightarrow O$ be the original injective map of
coefficient DVRs finite over $\mathbb Z_\ell$. Let $B_0,B$ be
local bridges for the same valuation ring and the same
\texttt{LocPrime}. Then for any integral matrix homomorphism $M_0$
(with no continuity assumption on $M_0$)
the extension, in the cocycle-span sense above, of its crystalline
local condition is exactly the crystalline condition of $f(M_0)$.

One inclusion follows by pushing a crystalline lift and its witness
through the coefficient map, then using SSC3 for the bridges.
To prove the reverse inclusion, note that $O$ is finite
torsion-free over $O_0$: finite generation over $\mathbb Z_\ell$
and SSC3 give finite generation over $O_0$. It is therefore free.
Choose an integral basis $b_a$ and fraction-field coordinate maps
$p_a:K\to K_0$. They satisfy
\[
 p_a(O)\subseteq O_0,\qquad
 p_a(f(y)x)=y p_a(x),\qquad x=\sum_a b_a f(p_a(x)).              \tag{SSC.7}
\]
They consequently induce maps $A_O^r\to A_{O_0}^r$ which
commute with the actions. Keep the source's literal witnesses for a
continuous quotient cocycle $c$ in its crystalline condition:
\[
 c(g)=[z(g)+M(g)v-v],
\]
where $z$ is a continuous crystalline $K$-valued cocycle and
$v\in K^r$. Apply each $p_a$ and SSC3 to $z$ and its period-ring
witness, and apply $p_a$ algebraically to $v$. The continuous
quotient coordinate map sends $c$ to a continuous cocycle, and
\[
 p_a(c(g))=[p_a(z(g))+M_0(g)p_a(v)-p_a(v)].
\]
Thus $p_a(c)$ has precisely the allowed original pair of witnesses
$(p_a(z),p_a(v))$. No continuity of the rational coboundary is
needed. Formula (SSC.7)
expresses the original quotient cocycle as a finite $O$-linear
combination of their pushes. This is exactly membership in
\texttt{extCond}, proving equality, including ramified coefficient
extensions.

Likewise any integral intertwiner $U\in\operatorname{GL}_r(O)$
induces inverse maps $[v]\mapsto[Uv]$ and $[v]\mapsto[U^{-1}v]$
on the quotient lattices and on continuous cochains. They commute
with restrictions and preserve inertia coboundaries, hence preserve
the unramified conditions. For the crystalline conditions use
$U\otimes1$ on the continuous lift's period-ring witness, and
apply $U$ separately to its rational boundary vector. Thus they give an actual
isomorphism of the full Selmer modules and preserve their lengths,
including infinite length.

For the source's $s,i$ and $\psi_L(i)$, the latter character is
trivial on every decomposition group above $\ell_i$. The local
representation there is exactly the coefficient extension of
$T_s$, and the first assertion applies. The source identities
$\ell_{\mathrm{twOf}(i)}=\ell_i$ and
$\mathrm{setS}_{\mathrm{twOf}(i)}=S_i$ identify all other local
conditions. Applying the integral conjugacy (SSC.6) yields the
full original equality
\[
        \texttt{sPiTop}(D,i,\rho_i,\pi)
             =\texttt{sTLoc}(D,i,s).                          \tag{SSC.8}
\]
The local field bridges and the integral lattices have remained
those in the original datum.
% END SSC4

\subsection{Arbitrary coefficient DVRs and the same character}

% BEGIN SSC5
\paragraph{SSC5: compatible coordinates in every torsion quotient.}
\label{claim:SSC5}
Let $g:O\hookrightarrow O'$ be any injective local homomorphism
of DVRs. There is no finiteness or completeness hypothesis here.
Choose a uniformizer $\varpi$ of $O$ and lift a vector-space basis
$(\bar b_a)_{a\in I}$ of $O'/g(\varpi)O'$ over $O/\varpi O$
to elements $b_a\in O'$. The index set may be infinite. For every
$0\ne d\in O$ these lifts give a basis of $O'/g(d)O'$ as an
$O/dO$-module.

It suffices to take $d=\varpi^n$. For generation, expand an
element modulo $\varpi$, subtract the finite expansion, divide
the remainder by $g(\varpi)$, and repeat $n$ times. Only finitely
many basis indices occur. For independence, reduce a putative
finite relation modulo $\varpi$. Every coefficient is divisible
by $\varpi$; cancel $g(\varpi)\ne0$ and repeat. Each coefficient
is divisible by $\varpi^n$. Multiplication by a unit proves the
assertion for general $d$.

Every element of $K'=\operatorname{Frac}(O')$ has a denominator
from $O$: $g(\varpi)$ has positive $O'$-valuation, so a large
power clears any denominator. Passing to the quotients by $O'$
and using the just proved generation and independence gives
\[
 \bigoplus_{a\in I} A_O^r \;\xrightarrow{\ \sim\ }\; A_{O'}^r,
 \qquad (x_a)\longmapsto\sum_a b_a\,\mathrm{push}(x_a).        \tag{SSC.9}
\]
Write $p_a$ for its $O$-linear coordinate maps. For matrices over
$O$ they commute with the action, and every vector has finite
coordinate support. These statements hold uniformly in every
denominator; a new basis has not been chosen at each torsion level.
% END SSC5

% BEGIN SSC6
\paragraph{SSC6: the full scalar-change theorem for normalized Selmer lengths.}
\label{claim:SSC6}
For the exact source's fixed $D,i,s,Q$, any injective
$f:O_{0,i}\to O$ and any injective local map $g:O\to O'$ of
DVRs, and any character $\psi:G\to O^\times$, one has
\[
 sT(D,i,s,Q,gf,g\psi)=sT(D,i,s,Q,f,\psi).                     \tag{SSC.10}
\]
Neither $O'$ nor $O$ is required to be finite over the other.

Apply (SSC.9) to $T_{s,f,\psi}$. Continuous cocycles with values
in a discrete module have finite image because $G$ is compact.
Thus each cocycle uses only finitely many of the coordinates $p_a$,
and the coordinate functions are continuous cocycles. The same
statement holds for a coboundary, using the finite support of its
defining vector. Consequently cohomology admits the corresponding
coordinate decomposition, compatible with restrictions to every
decomposition or inertia group. The latter groups are compact too.

It remains to check the exact crystalline-span condition, since a
coordinate map need not be $O'$-linear. A defining original cocycle
has finite image and hence is killed by one $0\ne d\in O$.
For any scalar $c\in O'$, expand $c$ modulo $g(d)O'$ using SSC5.
On all values of that cocycle, multiplication by $c$ is therefore
a finite combination of multiplication by the $b_a$ with
coefficients in $O$. Applying a coordinate to any finite sum
of such scalar multiples remains in the original defining
$O$-span. Pushes of allowed cocycles remain in the new span as
well. Intersecting with the cocycle modules and passing to
cohomology respects these maps. Inertia coboundaries are preserved
by direct equivariance, and the unrestricted local conditions add
no constraint. Thus, for the two Selmer modules $S,S'$, the natural
map is an isomorphism
\[
                     O'\otimes_O S\simeq S'.                 \tag{SSC.11}
\]
For clarity this assertion also holds for an arbitrary torsion
$S$ in the coordinate argument: any finite tensor sum is killed
by one $d$, and the bases of $O'/g(d)O'$ give both generation and
independence. This proves injectivity as well as surjectivity;
no exchange with an infinite product of local conditions is needed.
Each coordinate individually respects every local condition.

Let $e=v_{O'}(g(\varpi))>0$. Flatness of the torsion-free
$O$-module $O'$ and a composition series give
\[
 \operatorname{length}_{O'}(O'\otimes_O S)
          =e\operatorname{length}_O S.                        \tag{SSC.12}
\]
For infinite length, there are finite-length torsion submodules of
arbitrarily large length: otherwise maximal such length forces
the whole module to equal one finite submodule. Tensoring their
inclusions proves (SSC.12) in the infinite case. The denominator
$\operatorname{length}(O/\ell O)$ scales by the same $e$.
Cancellation proves (SSC.10). If the source's finite denominator
is zero, both finite normalized lengths are zero; if the numerator
has infinite length, both \texttt{nlen} values are infinity before
any division. This verifies the source conventions in all cases.
% END SSC6

% BEGIN SSC7
\paragraph{SSC7: the same-character relation produces a common DVR.}
\label{claim:SSC7}
Suppose $O_0,A,B$ are DVRs, $f_A:O_0\hookrightarrow A$ is
integral, $f_B:O_0\hookrightarrow B$ is local, and
$\psi_A:G\to A^\times$, $\psi_B:G\to B^\times$ satisfy the
original \texttt{SameChar}: every finite $O_0$-linear relation
among the first character values holds if and only if the
corresponding relation among the second values holds. There are a
DVR $O_1$, a map $f:O_0\to O_1$, injective local maps
$g_A:O_1\to A$, $g_B:O_1\to B$, and one character
$\psi:G\to O_1^\times$ whose images are exactly $\psi_A,\psi_B$,
with $g_Af=f_A$ and $g_Bf=f_B$.

Evaluate the group algebra $O_0[G]$ at the two characters. The
same-character condition says that the two kernels agree. Its
image $R\subset A$ therefore embeds in $B$. Let
$K=\operatorname{Frac}(R)$ and denote the induced field embeddings
by $j_A,j_B$. Put $O_1=K\cap A$, meaning the inverse image of $A$
under $j_A$. If $x\in O_1$, then $j_A(x)$ is integral over $O_0$
because $A$ is. A monic polynomial witnessing this also holds
for $j_B(x)$, so integral closedness of $B$ implies $j_B(x)\in B$.
This constructs both injections.

The injection into $A$ reflects divisibility: if $j_A(x)$ divides
$j_A(y)$, then $y/x\in K\cap A$. Hence divisibility in $O_1$
is total, and strict divisibility cannot descend indefinitely,
as its image cannot do so in a DVR. More concretely the valuation
of $A$ restricts to a subgroup of $\mathbb Z$ on $K^\times$;
$O_1$ is its nonnegative valuation ring. It is not a field:
a uniformizer of $O_0$ stays a nonzero nonunit in $B$ by locality
of $f_B$. Thus the restricted value group is a nonzero discrete
group, making $O_1$ a DVR. The map to $A$ is local by divisibility
reflection. If a uniformizer of $O_1$ became a unit in $B$, then
so would the nonzero nonunit from $O_0$, by its DVR factorization.
The map to $B$ is local as well. Character values and their inverses
already lie in $R$, so they define the required units of $O_1$.
Every stated equality follows from evaluation on group-algebra
generators. This proves the whole original common-DVR theorem.
% END SSC7

% BEGIN SSC8
\paragraph{SSC8: the actual same twist preserves Selmer length and local triviality.}
\label{claim:SSC8}
Apply SSC7 with $A=O_i$, $B=O_{\mathrm{twOf}(i),k}$,
$f_A=f_i$ and $f_B=\texttt{toOk}$. These maps have the required
properties. The first is integral because its target is finite
over $\mathbb Z_\ell$ and its two $\mathbb Z_\ell$ structures
agree by SSC3: a monic polynomial over $\mathbb Z_\ell$ is also
one over the source. For locality of the second map, use the
following elementary test. If $f:O\to B$ maps a nonzero natural
number $\ell$ to a nonunit, factor $\ell=u\varpi^n$. Then
$f(\varpi)$ is a nonunit. Every source nonunit is divisible by
$\varpi$, and hence cannot map to a unit. At the actual coefficient
rings $\ell$ is nonzero and remains a nonunit by integral finite
extension. This establishes locality.

The common character and common DVR of SSC7 give two applications
of SSC6, one to $g_A$ and one to $g_B$. Their middle normalized
length is identical. Thus the original same-character premise
gives the full equality
\[
 sTLoc(D,i,s)
      =sT(D,\mathrm{twOf}(i),s,Q_i,\texttt{toOk},\psi).        \tag{SSC.13}
\]
In addition, \texttt{SameChar} preserves equality of a character
value to $1$. Apply its defining relation to $[g]-[1]$ in the
group algebra. Therefore $\psi_L(i)=1$ on a decomposition group
implies $\psi=1$ there. In particular the transported character
is trivial at all places above $\ell$, exactly as required by the
following comparison. No assertion that the two original coefficient
rings are subrings of one another has been made.
% END SSC8

\subsection{From local crystalline classes to the signature Selmer module}

% BEGIN SSC9
\paragraph{SSC9: a cocycle-level inclusion into the original signature condition.}
\label{claim:SSC9}
For every $D,i,s,k,Q$ and $\psi:G\to O_{i,k}^\times$ trivial
on the decomposition groups above $\ell_i$, there is an injective
map from the source's \texttt{selT} to the signature Selmer module
$\texttt{Gen.selT}$. Both use the same $\psi$.

The coefficient map underlying this comparison is the elementary
isomorphism
\[
 K^r/O^r\longrightarrow(K/O)^r,
       [v]\longmapsto([v_1],\ldots,[v_r]).                    \tag{SSC.14}
\]
Its inverse takes any coordinate representatives. They differ
only by an integral vector. The formula
$[Mv]_a=\sum_b M_{ab}[v_b]$ shows equivariance. Inserting the
definitions of $T_{s,f,\psi}$, \texttt{chiS} and \texttt{Trep}
shows that these are the literal two representations in the source.

Here is the local verification, including the lift which its
definition requires. Let $c$ be an original quotient cocycle in
the crystalline condition. Its class is represented by the
reduction of a crystalline rational cocycle $z$, up to a
coboundary. Absorb the coboundary in $z$. For each $g$, choose a
representative $u_g\in K_0^r$ of $c(g)$. Under the local bridge,
the difference between $u_g$ and $z(g)$ is an integral vector.
Subtract its preimage in $O_0^r$ from $u_g$; call the result $y_g$.
Then $[y_g]=c(g)$ and the bridge sends $y_g$ exactly to $z(g)$.
This construction makes no assertion that an arbitrary choice
of representatives is continuous: continuity belongs to $c$ and
$z$, which are the objects required in the definition.

Push the $y_g$ to the new fraction field. To verify the source's
explicit fraction witness, write any coordinate $y=n/d$ with
$n,d\in O_0$, $d\ne0$. The identities
$dy=n$ survive both the bridge and the fraction-field map. These
are precisely the two denominator equations in the definition of
\texttt{Gen.bk}. The crystalline cocycle there is $z$, whose
representation is \texttt{Gen.repV}: both are the specified
inverse-character twist of $\rho_s$. Local triviality of $\psi$
ensures the coefficient push respects this local action.
Equation (SSC.14) now puts each defining original cocycle into
\texttt{Gen.bk}. Linear spans and the passage to cohomology preserve
this inclusion.

At the unrestricted places there is nothing to check; at all other
places (SSC.14) preserves inertia coboundaries. The original sets
$\mathrm{setS}$ and $\mathrm{Spl}$ coincide. Thus (SSC.14) induces
the claimed map of the global Selmer modules. It is injective
because its inverse is equivariant on the entire coefficient module
and hence on $H^1$. This proves the source's whole
\texttt{length\_le} comparison with no finite-length assumption.
% END SSC9

% BEGIN SSC10
\paragraph{SSC10: finite and infinite normalized lengths.}
\label{claim:SSC10}
Under the hypotheses of SSC9, write $S$ for the first Selmer module
and $S_{\mathrm{sig}}$ for the second. Then
\[
 \begin{split}
 \texttt{sLen}(S_{\mathrm{sig}})<\infty
      &\ \Longrightarrow\ n_O(S)<\infty,\\
 n_O(S)=q\in\mathbb Q
      &\ \Longrightarrow\ \operatorname{ofReal}(q)
                    \le\texttt{sLen}(S_{\mathrm{sig}}).
 \end{split}                                                   \tag{SSC.15}
\]
The injection of SSC9 gives the inequality of module lengths. If
the first length were infinite, the second would be infinite too.
Its extended nonnegative normalized length would then be infinite,
because the denominator is a finite natural number. This proves
the first assertion. For finite first length, division of the
length inequality by the common positive denominator proves the
second. This is the actual coefficient situation. At the broader
scope of the source's elementary \texttt{nlen\_le\_lenN} lemma,
the denominator may be zero. The finite rational normalization
on the left is then zero, so the second assertion is automatic;
the infinity argument for the first remains valid. This covers
both original conventions rather than silently identifying them.

Combining (SSC.8), (SSC.13), and (SSC.15) gives the full source
\texttt{sT\_le\_sLen} implication and its required application to
the actual point Selmer module. If that module is infinite, a finite
signature length is impossible. If it is finite, its rational
\texttt{sPi} value is exactly the rational value in (SSC.1).
% END SSC10

% BEGIN SSC11
\paragraph{SSC11: one attached signature proves the full Selmer core.}
\label{claim:SSC11}
Assume the source's \texttt{Cheb}, \texttt{TriplesGood},
\texttt{GoodPtRow}, and \texttt{ResAbsIrr}. At a point
$\pi\in Z_i^{\mathrm{good}}$, let $\rho_i(\pi)$ satisfy
\texttt{IsRhoPt}, and let the target character satisfy the full
original \texttt{SameChar} relation with $\psi_L(i)$ at the
specified coefficient maps. There is a signature $s$ which is
good, has the prescribed side, is attached to this actual point,
and has both conclusions (SSC.15) for its point Selmer module.

Choose the row $\sigma,S_g$ supplied by \texttt{GoodPtRow}.
Apply \texttt{TriplesGood} to that same row and parameter list.
This gives one $s$, its goodness, and its \texttt{IsTriple}.
Residual absolute irreducibility applies to the same
$\rho_i(\pi)$. SSC2 and SSC4 give (SSC.8); SSC8 gives
(SSC.13) and triviality of the target character above $\ell$;
SSC9--SSC10 give the two length conclusions. The triple's side
field gives the required side, and the collected row, parameter,
point and triple data give the original \texttt{AttachedIn}.
All five output fields therefore use the same $s$. This proves
the whole original \texttt{sel\_at} theorem, including arbitrary
target index $j$, its equality with $\mathrm{twOf}(i)$, and the
specified coefficient map transported by that equality.

Now take the original $C:\texttt{NChoice}$ and suppose
\texttt{NChTie} and \texttt{SameTwist}. At every $i$, the latter
provides the identical coefficient map after the prescribed
index identification and the character
\[
 \psi_G(C,w)=(C.\psi'(w))\circ\operatorname{art},
       \qquad w=D.\mathrm{nd}.w(i).                           \tag{SSC.16}
\]
Use the preceding construction with $k=C.k'(w)$. The tie identifies
$Q_i=C.Q(w)$ and the side with $C.\mathrm{sd}(w)$.
It gives every clause of \texttt{SelCore}, including
\texttt{AttachedIn}, for every actual family of \texttt{IsRhoPt}
representations. In particular it applies to the \texttt{PtFam}
constructed in DFD7, because $Z_i^{\mathrm{good}}\subseteq Z_i$.
% END SSC11

\subsection{Tame conductors, residual irreducibility, and the final join}

% BEGIN SSC12
\paragraph{SSC12: normalized valuations agree in the two coefficient realizations.}
\label{claim:SSC12}
Let a field $L$ with valuation $w$ embed into fraction fields of
DVRs $R_1,R_2$, and suppose both embeddings identify the condition
$w(x)\ge0$ with integrality. Let $\ell$ be nonzero and a nonunit
in both DVRs. If $x\in L$ has integral images $a,b$, then
\[
             \frac{v_1(a)}{v_1(\ell)}
                  =\frac{v_2(b)}{v_2(\ell)}.                 \tag{SSC.17}
\]
At $x=0$ use the source's $\mathrm{toNat}(\infty)=0$ on both
sides. Otherwise, for all natural $p,q$,
\[
 \ell^q\mid a^p
   \ \Longleftrightarrow\ x^p/\ell^q\text{ is integral}
   \ \Longleftrightarrow\ \ell^q\mid b^p.
\]
Write $A_j=v_j(\ell)>0$, $B_1=v_1(a)$, $B_2=v_2(b)$.
This says $qA_1\le pB_1$ if and only if $qA_2\le pB_2$.
Use $(p,q)=(A_1,B_1)$ and $(A_2,B_2)$ to obtain both
inequalities between $B_1A_2$ and $B_2A_1$, proving (SSC.17).

Also
\[
 \operatorname{lenUnit}_{\ell}(R/(x))
       = nval_{R,\ell}(x),\qquad
 nval_{R,\ell}(ux)=nval_{R,\ell}(x)\ (u\in R^\times).         \tag{SSC.18}
\]
These follow from the uniformizer factorization and the length of
$R/\varpi^nR$. For $x=0$ both sides of the first equality use the
same zero convention; it is not a claim that $R$ has finite length.
The same valuation argument proves invariance under local
coefficient extensions, by canceling the common ramification index.
% END SSC12

% BEGIN SSC13
\paragraph{SSC13: tame generators and inversion of the character.}
\label{claim:SSC13}
For the actual $v\in Q_i$, the finite-order character $\psi_L(i)$
is trivial on wild inertia. Here is an elementary verification of
the exact input needed. A wild automorphism $\tau$ satisfies
$|\tau(y)/y-1|_v<1$ for every nonzero $y$, and so do all its
nonnegative powers, by expanding $(1+a)(1+b)-1$.
If $\tau^k x=x$ with $|k|_v=1$, then $\tau x=x$. Otherwise put
$d=\tau x-x\ne0$; telescoping gives
\[
 k=-\sum_{j=0}^{k-1}\left(\frac{\tau^j(d)}d-1\right),
\]
contradicting $|k|_v=1$. The open kernel of $\psi_L(i)$ contains
the subgroup fixing some finite extension. A power $\tau^N$
fixes that extension. Write $N=\ell^a m$ with $\ell\nmid m$.
Since $v\nmid\ell$, the preceding observation shows that
$\tau^m$ already fixes it. Thus $\psi_L(i)(\tau)^m=1$.
This value also has $\ell$-power order by the original character
field, so it equals $1$.

Let $\gamma=\gamma_{i,v}$ be the exact supplied tame generator.
The inverse image of $\langle\psi_L(i)(\gamma)\rangle$ is an open,
and hence closed, subgroup. It contains $\gamma$ and wild inertia.
The supplied tame-generator density therefore puts all inertia
inside this inverse image. Thus each inertia value is an integral
power of $\psi_L(i)(\gamma)$. The shared character $\psi_0$
and both injective coefficient realizations in \texttt{NChTie}
transport the identical equality of powers to $\psi_G(C,w)$.
This proves the original full \texttt{TameCyc} conclusion.

For any unit $u$ and integer $m$ with $u^m\ne1$, $u-1$ divides
$u^m-1$: use the geometric sum for $m\ge0$ and multiply by
$-u^m$ for negative $m$. Hence
$nval(u-1)\le nval(u^m-1)$. The source's tame value is the
infimum of these normalized valuations over nontrivial inertia
values. If the character value at $\gamma$ is nontrivial, it
belongs to that set and gives a lower bound for every element;
thus
\[
 nval(\psi_G(\gamma)-1)\le\operatorname{tameVal}(\psi_G).
                                                                    \tag{SSC.19}
\]
If it is trivial, all inertia values are trivial and the defining
set is empty; the source uses $\inf\varnothing=0$ in $\mathbb R$,
so the same inequality holds. Finally
$u^{-1}-1=-u^{-1}(u-1)$ implies
\[
           \operatorname{tameVal}(\psi^{-1})
                     =\operatorname{tameVal}(\psi).           \tag{SSC.20}
\]
This includes the trivial-character case and retains the actual
generator and coefficient embeddings.
% END SSC13

% BEGIN SSC14
\paragraph{SSC14: the full conductor-root bound at the tied auxiliary places.}
\label{claim:SSC14}
Assume \texttt{NChTie} for $D,C$. For every $i$ and every
$v\in Q_i$,
\[
 \operatorname{lenUnit}_{\ell_i}
        \bigl(O_i/(\zeta_{i,v}-1)\bigr)
  \le \operatorname{tameVal}_{\ell_i}
          (\psi_W(C,w),v),\qquad w=D.\mathrm{nd}.w(i).         \tag{SSC.21}
\]
Indeed $Q_i$ is disjoint from $S_i$ and from the primes above
$\ell_i$. The first twist character is unramified at $v$, since
all its possible ramification places belong to $S_i$. Its value
on $\gamma_{i,v}$ is therefore $1$. The exact character product
in the datum gives
$\zeta_{i,v}=\psi_L(i)(\gamma_{i,v})$.

The two realizations of $\psi_0$ in \texttt{NChTie} identify
$\psi_L(i)(\gamma)-1$ and the image of
$\psi_G(C,w)(\gamma)-1$ in $C.O_e(w)$ as images of one element
of the original number field. Both embeddings induce the same
specified place. The prime $\ell_i$ is nonzero and a nonunit in
all these rings: first in the original coefficient rings, and then
in the finite integral extensions defining $O_{i,k}$ and $C.O_e$.
SSC12 thus identifies their normalized valuations, and its last
assertion removes the extra coefficient extension. By definition
\[
 \psi_W(C,w)=\operatorname{psiAt}(C.\mathrm{sd}(w),C.\psi'(w))
                         \circ\operatorname{art}.
\]
It is either $\psi_G$ or $\psi_G^{-1}$. Apply (SSC.19)--(SSC.20)
and (SSC.18) to obtain (SSC.21). Summing it and using the literal
tie $Q_i=C.Q(w)$ gives the full second clause of
\texttt{SignatureChoiceCondition}. The zero and empty-set cases
have already been included; no strict positivity of a conductor
root is assumed.
% END SSC14

% BEGIN SSC15
\paragraph{SSC15: the original residual hypothesis covers every chosen point representation.}
\label{claim:SSC15}
The source's \texttt{ResidualIrreducibleRankR} implies its exact
\texttt{ResAbsIrr}. To see this, fix $\pi\in Z_i^{\mathrm{good}}$
and any \texttt{IsRhoPt} realization $\rho$. Its actual rational
block decomposition gives the characteristic-polynomial identity
\[
 \operatorname{charpoly}(R(g))
    =\operatorname{charpoly}(\rho(g))
           \prod_{c=1}^{t}(X-\chi_{i,c}(g)).                   \tag{SSC.22}
\]
Both sides have coefficients in $O_i$ and the fraction-field map
is injective, so (SSC.22) is an integral equality. Reduce it modulo
the maximal ideal and use the specified residual Hecke-polynomial
tie at the permitted Frobenius elements. Together with the open
kernel of the continuous residual representation, these are
exactly the fields of \texttt{IsRhoBar}. The good point also gives
the source's \texttt{HG} premise. The original residual hypothesis
therefore makes the representation irreducible over the given
algebraically closed residual coefficient field.

This is absolute irreducibility over the original residue field.
Explicitly embed its algebraic closure into that given field.
A nonzero invariant subspace extends to a nonzero invariant
subspace there. A proper subspace remains proper: choose a
nonzero linear functional killing it and extend that functional
by the field embedding. Thus any nonzero proper invariant subspace
would contradict the stated irreducibility. The rank is positive
by the original ranks record. This proves \texttt{AbsIrrSub} on the
full space and hence \texttt{ResAbsIrr}, uniformly for every chosen
\texttt{IsRhoPt} lattice, not only a preferred representative.
% END SSC15

% BEGIN SSC16
\paragraph{SSC16: the whole original Selmer-signature supplier.}
\label{claim:SSC16}
The following is now proved at its full original scope:
\[
\begin{gathered}
 \texttt{Cheb},\quad \texttt{PtFam}(D,R,\rho),\quad
 \texttt{NChTie}(D,C),\\
 \texttt{TriplesGood}(D),\quad\texttt{SameTwist}(D,C),\quad
 \texttt{GoodPtRow}(D),\quad\texttt{ResAbsIrr}(D)\\
 \Longrightarrow\quad
            \texttt{SignatureChoiceCondition}(\rho,C).        \tag{SSC.23}
\end{gathered}
\]
Apply SSC11 to each $\pi\in Z_i^{\mathrm{good}}$ and use the
\texttt{PtFam} field on that point. Its signature has side
$\mathrm{side}(s)=C.\mathrm{sd}(w)$. By the two-case definition
of \texttt{psiAt},
\[
 \operatorname{psiAt}(C.\mathrm{sd}(w),\psi_W(C,w))
                        =\psi_G(C,w).                         \tag{SSC.24}
\]
For the true side nothing is inverted; for the false side the
same character is inverted twice. The side equality is essential:
two unrelated Boolean choices could not be canceled. Therefore
both the finiteness implication and the normalized-length
inequality of SSC11 are exactly the first clause of the required
condition, with its specified signature character. SSC14 supplies
its second clause. These statements hold at every location and
therefore at every original \texttt{TwLoc}; no enlargement of the
location set is required by the theorem. This proves the entire
original \texttt{JointSix.main} supplier, retaining the supplied
\texttt{Cheb} premise and every other original input.

With \texttt{TwGoodAt} on the same $C$ and $I$, DFD11 now gives
$\texttt{KillAt}(\rho)$ directly. The budget remains exactly
$1/\ell^2$ and the local term is the same tame term in (SSC.21).
This conclusion is about the actual chosen representation family.
It neither replaces the family nor chooses new twists after
the length estimates.
% END SSC16

% BEGIN SSC17
\paragraph{SSC17: the exact family setup now needs only the coherent Wach supplier.}
\label{claim:SSC17}
In DFD15 retain the original family, constants, \texttt{TwGood},
\texttt{FactsF}, \texttt{Cheb}, and just its first extra supplier:
for every relevant $D$, actual compatible
$HW,F_W,\texttt{CmpOf}$ on $D.\mathrm{loc},D.\mathrm{locA}$.
Then the full original output of \texttt{exists\_setup\_large}
follows. Its separate Selmer-signature supplier is no longer an
extra premise.

Choose once the maximum of the parameter, small-prime, and
\texttt{TwGood} cutoffs in DFD13--DFD15. For any later
$x,h_x,i$, retain the specified $D,L,C_D=\texttt{nch}$ and twist
index $I$. Choose the coherent $HW,F_W,\mathrm{Cmp}$ using the
remaining supplier. DFD9 constructs the actual canonical
$X$ with $X.I=I$, error function zero and exceptional-set function
constantly the arbitrary given $\mathrm{exc}$. In particular
$X.\mathrm{exc}=\mathrm{exc}$, and its pair $(X.R,X.\rho)$ has
the original \texttt{PtFam}.

The member of \texttt{FactsF} supplies \texttt{NChTie},
\texttt{TriplesGood}, and \texttt{SameTwist}. DFD10 supplies
\texttt{GoodPtRow} from that member's actual row conclusions.
Its residual rank-$r$ hypothesis gives \texttt{ResAbsIrr} by
SSC15. SSC16 applies to this precise $X.R,X.\rho,C_D$, so DFD11
gives the killing field. The facts, canonical, two parameter,
small-scale, split-bound and supplied \texttt{ShaFinOdd} fields
are the other seven fields already assembled in DFD15.
The same parameter implication gives \texttt{AdmJ}.

The cutoff is chosen before $x,h_x,i$, before the
\texttt{ShaFinOdd} premise, and before every finite exceptional
set. All objects in the output are the ones just constructed.
This proves the exact conditional setup statement with only the
coherent local comparison supplier still displayed. DFD12 proves
its finite-coefficient Bhatt--Scholze reduction, but construction
of the whole coherent Wach system remains. The later local row
bounds and the complete common \texttt{ConclF} family also remain;
\texttt{GLOBAL\_DENOMINATOR} is consequently still open.
% END SSC17

\subsection{Source alignment and limits of the result}

The pinned source is commit
\texttt{7cbf7a0c32fca1badf05931879a37c3f8a2b03b4}.
The accompanying source manifest records the exact files, ranges,
bytes and hashes. In the \texttt{InertiaEigenvalues} proofs the
main endpoints are the trace and twisted-conjugacy theorems,
\texttt{SPiTopEqSTLoc}, \texttt{IsLocalHomOfNatCast},
\texttt{NlenLeLenN}, \texttt{SameCharEqOne},
\texttt{ConductorRoot\allowbreak SubOneLength\allowbreak LeTameValue},
\texttt{ResidualAbsIrred\allowbreak GoodPointOf\allowbreak ResidualIrred}, and
\texttt{PsiAtPsiAt}. The foundation files for integral trace
conjugacy, common coefficient rings and arbitrary injective local
coefficient changes are reconstructed in SSC1 and SSC5--SSC8;
their names are not used as substitute proofs.

Classical algebraic foundations used here are the
finite-dimensional density theorem, field-embedding extension,
basic DVR structure and module length, and the usual compactness
and valuation properties of absolute Galois groups and finite
$\ell$-adic fields. The original local period-ring and bridge
structures, original tame-generator density, and original
\texttt{Cheb} are retained where the source assumes them.
No new automorphy, crystalline existence, or Wach comparison
theorem is inserted among these foundations.
